\subsection{Konstruktion von Tests}
% Slide 401 (P24)
$\cX_i$ diskret oder gemeinsam stetig unter $\P_{\vartheta_0}$ und $\P_{\vartheta_A}$. Sei $L(x_1, \ldots, x_n; \vartheta)$ die Likelihood-Funktion

\shortdefinition[Likelihood-Quotient] Für $\vartheta_0 \in \Theta_0$, $\vartheta_A \in \Theta_A$ und $x_i \in \R$:
$R(x_1, \ldots, x_n; \vartheta_0, \vartheta_A) = \frac{L(x_1, \ldots, x_n; \vartheta_A)}{L(x_1, \ldots, x_n; \vartheta_0)}$.\\
$R(x_1, \ldots; \vartheta_0, \vartheta_A) = +\8$ wenn $L(x_1, \ldots; \vartheta_0) = 0$.
Wenn $R$ gross ist, sind Beobachtungen in $\P_{\vartheta_A}$ deutlich wahrscheinlicher als in $\P_{\vartheta_0}$.

% TODO: More places with highlighting
\shortdefinition[Likelihood-Quotienten-Test] Sei $c \geq 0$. Der LQT mit param $c$ ist Test $(T, K)$ mit $T = R(\cX_1, \ldots, \cX_n; \vartheta_0, \vartheta_A)$ 
und $K = (c, \8)$. $H_0$ wird verworfen wenn $R$ gross wird.

\shorttheorem[Neyman-Pearson] Seien $\Theta_0, \Theta_A$ \textit{einfach} und $(T, K)$ ein LQT mit $c$ und sig.-niv. $\alpha^* := \P_{\vartheta_0}[T \in K]$.
Anderer Test $(T', K')$ mit $\P_{\vartheta_0}[T' \in K'] =: \alpha \leq \alpha^*$, so gilt auch $\P_{\vartheta_A}[T' \in K'] \leq \P_{\vartheta_A}[T \in K]$.

\inlineintuition Jeder Test mit kleinerem Signifikanzniveau hat auch kleinere Macht

\shortdefinition[Verallgemeinerte LQ] ist gegeben durch
\[
    R(x_1, \ldots, x_n) = \frac{\sup_{\vartheta \in \Theta_A} L(x_1, \ldots, x_n; \vartheta)}{\sup_{\vartheta \in \Theta_0} L(x_1, \ldots, x_n; \vartheta)}
\]
